\documentclass[11pt]{article}
\usepackage[margin=1in]{geometry}
\usepackage{amsmath,amssymb,amsthm}
\usepackage{booktabs}
\usepackage{hyperref}

\newtheorem{theorem}{Theorem}
\newtheorem{conjecture}[theorem]{Conjecture}
\newtheorem{claim}[theorem]{Claim}
\theoremstyle{definition}
\newtheorem{reading}[theorem]{Reading}
\newtheorem{counterexample}[theorem]{Counterexample}

\title{Disproof of TLMC Conjecture 00000001211:\\ P-Positions of the Young Poset Game Are Not the Perfect Squares}
\author{}
\date{}

\begin{document}
\maketitle

\section*{Conjecture and verdict}

\begin{conjecture}[TLMC 00000001211, verbatim]
Definition: A poset game (delete an element and everything above it). Conjecture: The
P-positions of the poset game on Young's lattice have an explicit characterization via
perfect square partitions (the Young poset game).
\end{conjecture}

\medskip
\noindent\textbf{Verdict: disproved.} The phrase ``poset game on Young's lattice'' admits two
natural readings; under \emph{both} the perfect-square characterization fails, and every number
in this note is certified twice: by a Lean development that uses \emph{no axioms at all} (no
\texttt{sorry}, no \texttt{Classical.choice}, no \texttt{Quot.sound}, not even
\texttt{propext}; see \texttt{lean4/Check.lean}) and by an independent Python enumeration over
all $252$ Young diagrams fitting in the $5\times 5$ box (\texttt{reproduce.py}).

\section*{Setup}

A \emph{Young diagram} is a partition $\lambda=(\lambda_1\ge\lambda_2\ge\dots\ge 0)$, drawn as
rows of cells, row $i$ (1-indexed, from the top) having $\lambda_i$ cells. A move
\emph{chomps} a cell $(i,j)$: the cell together with every cell $(r,c)$ with $r\ge i$, $c\ge j$
is removed. Under normal play, a position is a \emph{P-position} if the player to move loses.
A \emph{perfect-square partition} is $(k,k,\dots,k)$ with $k$ rows of length $k$
(a $k\times k$ square).

\begin{reading}[Literal poset game]
Every cell of the current diagram may be chomped; the player who cannot move (facing $\emptyset$)
loses.
\end{reading}

Since the corner cell $(1,1)$ belongs to every nonempty diagram and its up-set is the whole
diagram, chomping $(1,1)$ clears the board in one move.

\begin{theorem}[Reading 1]
The P-positions of the literal Young poset game are exactly $\{\emptyset\}$. In particular the
perfect square $(2,2)$ is an $N$-position.
\end{theorem}

\begin{proof}[Lean certificate]
\texttt{isPLit\_nil : isPLit [] = true} and \texttt{isPLit\_cons : 1 <= n -> isPLit (n :: rest)
= false}, i.e.\ no nonempty position is a P-position; combined they pin $P=\{\emptyset\}$.
The key lemma \texttt{chompAt\_11} states that chomping $(1,1)$ always yields the empty
diagram; the witness move then makes every nonempty position an $N$-position.
\end{proof}

\begin{counterexample}[Reading 1]
$(2,2)$ is a perfect-square partition but an $N$-position (indeed so is every nonempty
diagram). If one insists on regarding $\emptyset$ as the $0\times0$ square, the inclusion
``$P\subseteq$ squares'' holds trivially while ``squares $\subseteq P$'' fails; either way the
explicit characterization is false.
\end{counterexample}

\begin{reading}[Chomp with a poison corner]
The corner $(1,1)$ is poison: taking it loses. Equivalently---since a rational player never
takes poison voluntarily---forbid that move; a player facing the lone poison cell $(1)$ has no
winning move and must take the poison, so $(1)$ is a P-position, exactly as in Chomp.
\end{reading}

\begin{theorem}[Reading 2]
The characterization fails in both directions: $(2,1)$ is a P-position that is not a square,
and the perfect squares $(2,2)$ and $(3,3,3)$ are $N$-positions.
\end{theorem}

\begin{proof}[Lean certificate]
\texttt{chomp\_counterexample} bundles the \texttt{decide}-verified evaluations
\texttt{isPChomp [2,1] = true}, \texttt{isSquare [2,1] = false},
\texttt{isSquare [2,2] = true}, \texttt{isPChomp [2,2] = false},
\texttt{isSquare [3,3,3] = true}, \texttt{isNChomp [3,3,3] = true}.
\end{proof}

\section*{The 5\texorpdfstring{$\times$}{x}5 enumeration}

There are $\binom{10}{5}=252$ partitions fitting in the $5\times5$ box. The literal reading
yields the single P-position $\emptyset$. The Chomp reading yields exactly $24$ P-positions:
\[
\begin{gathered}
(), (1), (2,1), (2,2,1), (3,1,1), (3,2), (2,2,2,1), (4,1,1,1), (4,3), (3,3,1,1), (4,2,2),\\
(2,2,2,2,1), (4,2,1,1,1), (5,1,1,1,1), (5,2,1,1), (5,4), (3,3,2,1,1), (5,3,2), (3,3,3,2,2),\\
(4,4,3,1,1), (5,3,3,2), (5,5,3), (4,4,2,2,2), (5,5,2,2).
\end{gathered}
\]
The only square among them is the degenerate $(1)$; the nondegenerate squares
$(2,2),(3,3,3),(4,4,4,4),(5,5,5,5,5)$ are all $N$-positions (as is the $2\times3$ rectangle
$(3,3)$).

The $3\times3$ corner of the table, under both readings:

\begin{center}
\begin{tabular}{lcc}
\toprule
position & Reading 1 (literal) & Reading 2 (Chomp) \\
\midrule
$\emptyset$    & P & --- (not reached) \\
$(1)$          & N & P (poison) \\
$(2)$          & N & N \\
$(1,1)$        & N & N \\
$(2,1)$        & N & \textbf{P} (not a square) \\
$(2,2)$        & N & \textbf{N} (a square) \\
\bottomrule
\end{tabular}
\end{center}

For instance, under Reading 2 the moves of $(2,1)$ chomping cells $\ne(1,1)$ are
$(1,1)$ and $(2)$, both $N$ (each has the move to $(1)$), so $(2,1)$ is $P$; then $(2,2)$ has
the move $(2,2)\to(2,1)$ to a P-position and is $N$.

\section*{Machine-checked certificate}

\texttt{lean4/Main.lean} encodes diagrams as \texttt{List Nat} (non-increasing row lengths)
and defines the chomp move \texttt{chompAt}, the cell enumeration \texttt{anyCell}/\texttt{cols},
and two fuel-based win oracles \texttt{winLit}/\texttt{winChomp} with the $N/P$ predicates
\texttt{isNLit}/\texttt{isPLit}/\texttt{isNChomp}/\texttt{isPChomp} at fuel
$\mathrm{total}(\lambda)+1$. The fuel is sufficient: every move enumerated is a genuine cell
(column $j$ of a row of length $n$ has $1\le j\le n$), hence row $i$ shrinks from $len$ to
$\min(len, j-1)\le len-1$ and every move removes at least one cell, so the game depth from a
diagram with $t$ cells is at most $t$. Chomped rows that become empty end the construction
immediately (for genuine cells of canonical partitions this can only happen when $j=1$, which
empties all lower rows too), so results stay canonical partitions; this also keeps every
definition free of \texttt{propext}.

\texttt{lean4/Check.lean} prints \texttt{\#print axioms} for all $18$ theorems; every line
reads \emph{``does not depend on any axioms''}. \texttt{reproduce.py} re-derives the
enumeration independently and asserts every figure quoted above.

\section*{Conclusion}

Conjecture 00000001211 is false under both readings of the game: the P-positions of the Young
poset game are $\{\emptyset\}$ (literal reading) or a $24$-element set inside the $5\times5$
box that neither contains nor is contained in the set of perfect-square partitions (Chomp
reading). No perfect-square characterization of the conjectured form exists.

\end{document}
